\scanpage{60}
c'est le groupe des autom du carré combinatoire
\deleted{\unsure{permutant|permute} entre elles de} qu'elles forment
\inserted{($\simeq\mathfrak S_{5}$)}, \deleted{\unsure{qui} se forment
\unsure{orientations}}. L'ext. de $\mathbb Z/2\mathbb Z$ par $\operatorname{Aut}^+(I)$
est l'image inverse du sous-groupe \unsure{d'|de}
ordre 2 \inserted{(qui fixe \unsure{le|les} \deleted{deux} \unsure{côté|côtés} du carré correspondant
aux choix des deux orientations (mais qui peuvent échanger ses deux extrémités).}
\deleted{\unsure{partie du cr} \unsure{\illegible{2}|de l'antipodisme} de carré}
\deleted{de $\{\pm1\}$ de $\mathbb F_5^*$, c'est le groupe
des autom du système, respectant la
paire d'orientations (sans nec$^t$ fixer
la comp. conn. corresp.). L'ext.}
L'ext. de $\mathbb Z/4\mathbb Z$ par $\operatorname{Aut}^+(I)$ s'obtient en prenant
l'image inverse du groupe des rotations
du carré. L'intersection de ces deux
sous-groupes de $G$ est $\operatorname{Aut}^+(I)$
(\unsure{se projette} dans le groupe des autom du
carré). \deleted{\unsure{L'é}}Le groupe engendré est formé
des produits, et c'est $G$ tout entier \ldots

\uuline{Généralisation} aux cartes de congruences.

La carte $\unsure{\mathrm{Cong}|\mathrm{Congr}}_{n}$ étant réalisée avec \deleted{des}\inserted{l'ens des sommets $S$ formé}
\deleted{sommets qui sont des} paires $\{a,-a\}$
d'éléments d'un $\mathbb Z/n\mathbb Z$-module libre \inserted{symplectique $E$} de
rang $2$, formant base, on peut y
% slanted
\margin{Supposer $n$ premier i.e. $\mathbb Z/n\mathbb Z$ un corps, ici ??}
associer $a\otimes a\in\operatorname{End}(E)=M$, \inserted{(alg. Az.)} ce sont
des éléments du cône isotrope $C$ de $M^{\unsure{\natural|\sharp|0}}_{\unsure{s|8}}$
(éléments de trace nulle) \unsure{i.e.|ipso} \unsure{a priori|donc ici|formé} des
nilpotents. On trouve ainsi une réalisation
% archive sheet number: 59
